% Part: first-order-logic
% Chapter: beyond
% Section: other-logics

\documentclass[../../../include/open-logic-section]{subfiles}

\begin{document}

\olfileid{fol}{byd}{oth}

\olsection{অন্যান্য যুক্তিবিদ্যা}

এতক্ষণে নিশ্চয়ই বোঝা গেছে, নতুন একটি যুক্তিবিদ্যা নকশা করা কঠিন
নয়। তুমিও নিজস্ব একটি বাক্যতত্ত্ব তৈরি করতে, একটি অবরোহী ব্যবস্থা
উদ্ভাবন করতে এবং তার উপযোগী একটি অর্থতত্ত্ব গড়ে নিতে পারো।
!!{derivation} ব্যবস্থাটি অর্থতত্ত্বের জন্য সম্পূর্ণ করতে চাইলে অবশ্য
একটু বুদ্ধি খাটাতে হতে পারে, আর তোমার যুক্তিবিদ্যাটি সত্যিই আকর্ষণীয়
বলে বৃহত্তর জগৎকে বোঝাতে কিছু পরিশ্রমও লাগতে পারে। বদলে নিজের
যুক্তিবিদ্যার গাণিতিক ও গণনামূলক ধর্মগুলি অনুসন্ধান করে নির্মল আনন্দে
অনেক ঘণ্টা কাটাতে পারবে।

সাম্প্রতিক দশকগুলিতে আনুষ্ঠানিক যুক্তিবিদ্যার যেন বিস্ফোরণ ঘটেছে।
অস্পষ্ট ধর্ম নিয়ে যুক্তিবিচারের মডেল হিসেবে ফাজি যুক্তিবিদ্যা নির্মিত।
অনিশ্চয়তা নিয়ে যুক্তিবিচারের মডেল হিসেবে সম্ভাবনামূলক যুক্তিবিদ্যা
নির্মিত। ডিফল্ট যুক্তিবিদ্যা ও অ-একঘেয়ে যুক্তিবিদ্যা প্রত্যাহারযোগ্য
যুক্তিবিচারের মডেল হিসেবে নির্মিত; অর্থাৎ, নতুন তথ্য এলে পরে উল্টে যেতে
পারে এমন “যুক্তিসঙ্গত” অনুমানের মডেল এগুলি। জ্ঞান নিয়ে যুক্তিবিচারের
মডেল হিসেবে জ্ঞানতাত্ত্বিক যুক্তিবিদ্যা, কারণগত সম্পর্ক নিয়ে
যুক্তিবিচারের মডেল হিসেবে কারণমূলক যুক্তিবিদ্যা, এমনকি নৈতিক ও
নীতিশাস্ত্রীয় কর্তব্য নিয়ে যুক্তিবিচারের মডেল হিসেবে “কর্তব্যগত”
যুক্তিবিদ্যাও আছে। এই ব্যবস্থাগুলি প্রবর্তনের প্রধান প্রেরণা দার্শনিক,
গাণিতিক না গণনামূলক, তার উপর নির্ভর করে এমন সৃষ্টিগুলির অধ্যয়ন
গাণিতিক যুক্তিবিদ্যা, দার্শনিক যুক্তিবিদ্যা, কৃত্রিম বুদ্ধিমত্তা,
জ্ঞানবিজ্ঞান বা অন্য কোনো শিরোনামের অধীনে পাওয়া যেতে পারে।

তালিকা ক্রমেই দীর্ঘ হয়, আর সম্ভাবনাগুলিকে মনে হয় অন্তহীন। মানুষের
সমস্ত বিচারবুদ্ধিকে গণনায় নামিয়ে আনার লাইবনিৎসের স্বপ্ন হয়তো কখনোই
আমরা অর্জন করব না---কিন্তু তাতে চেষ্টা থামানোর কারণ নেই।

\end{document}
